\chapter{Method and Evaluation Design}
\label{ch:eval}

% SOURCES OF TRUTH:
%   thesis-pack/09-method/EVALUATION_PROTOCOL.md  metric definitions and gates
%   thesis-pack/09-method/PREREGISTRATION.md      hypotheses, margins, k, models,
%                                                 and the 16 amendments
%   thesis-pack/00-start-here/RESULTS.md          wins on any executed number
%   thesis_doc/rubric-bibliography.md             the defence for each rubric choice
%
% RULES FOR THIS CHAPTER:
%  * Describe the design AS REGISTERED, in the present tense. Say where the
%    execution differed, in the past tense.
%  * The design's own failure (Gate 0, the null policy at the ceiling) belongs
%    to Chapter~\ref{ch:critique}. Here it gets a factual pointer only.
%  * Never write that the baseline "adapted", "recovered" or "replanned better".

This chapter fixes what the evaluation measured and how it decided. It is
written as the design was registered before the evaluation ran. Where the
execution differed from the plan, the difference is stated at the place where
it matters, with the amendment that records it. The results are reported in
Chapter~\ref{ch:results}. The criticism of the design, including the finding
that its main test could not resolve in either direction, is
Chapter~\ref{ch:critique}.

The thesis is a design-and-evaluation study. The method therefore has one job.
It has to make the comparison between the two systems fair enough that its
answer can be trusted, whichever way that answer goes. Every section below
answers one objection an examiner could raise: that the comparison was rigged,
that the baseline was a strawman, that the thresholds were picked to produce a
result, or that the instrument did not measure what it claims to measure.

% =====================================================================
\section{Study Design}
\label{sec:eval:overview}

The study compares two coordination systems on the same evacuation task. The
\emph{treatment} is the database-centric agentic system of
Chapter~\ref{ch:agentic}, where distributed LLM agents reason about the
situation and negotiate through the shared database. The \emph{control} is the
workflow baseline of Chapter~\ref{ch:workflow}, a single scripted Incident
Commander that follows written coordination doctrine.

There is one independent variable: the coordination architecture at the
command post. Everything else is held constant. It is held by reuse, not by
matching. The baseline engine is a subclass of the agentic engine, and 27 of
its attributes are the very same function objects that the live agents use,
not copies of them. Chapter~\ref{ch:validity} audits this in detail.

\paragraph{Pairing.} Runs are paired. A treatment run and a control run of the
same scenario receive byte-identical inputs through the same loader and write
to the same append-only \code{event_log}. A reference protocol for this kind of
comparison splits randomness into seven streams and asks that the six external
ones be matched across a pair. In this simulator six of the seven are constant
by construction. The map is a fixed set of real addresses, the flood is one
declared arrival time per zone, the head-count is declared per scenario,
residents have no behaviour model, and there is no sensor or message-latency
model. Only the policy stream is live. Matching the other six is therefore
trivial.

The live stream is controlled by a seed. Each run records
\begin{equation}
\texttt{policy\_seed} = \operatorname{sha256}\big(\texttt{"<scenario>|<scale>|<replicate>"}\big) \bmod 2^{31}.
\end{equation}
The key leaves out the system and the model on purpose. A treatment run and its
paired control therefore draw the same seed, and a test fails if either field
is ever added. The seed is passed to every model call. A seeded call is still
only \emph{near}-deterministic, because batching and cache state move the
sampling. The claim is therefore ``seeded and logged'', never ``reproducible
at the inference level''. This is the reason for replication.

\paragraph{Unit of analysis.} The unit is the run, not the disruption. Each
cell of system, scenario, model and time scale is repeated $k = 10$ times, with
replicate indices 1000 to 1009. Section~\ref{sec:eval:analysis} explains how
$k$ was set. The baseline runs no language model. For each scenario it is a
deterministic reference point, not a distribution. The variance of the agentic
side is reported as a cost of the architecture and is never averaged away.

% =====================================================================
\section{Hypotheses}
\label{sec:eval:hypotheses}

Four hypotheses were registered before any evaluation run. They are listed in
Table~\ref{tab:eval:hypotheses}. H1 and H2 are confirmatory and together they
answer RQ1. H3 and H4 are estimation questions and carry no significance test.

\begin{table}[htbp]
\centering
\caption{The registered hypotheses.}
\label{tab:eval:hypotheses}
\small
\begin{tabularx}{\textwidth}{l>{\raggedright\arraybackslash}Xll}
\toprule
 & Hypothesis & Arm & Form \\
\midrule
H1 & The agentic system is not worse than the baseline on survival, on disruptions the baseline has a written branch for. & Anticipated & Non-inferiority, $\delta_{NI} = 0.05$ \\
H2 & The agentic system is better than the baseline on disruptions the baseline could not have encoded. & Novel & Superiority, $\delta_m = 1$ level \\
H3 & An H2 advantage holds across a range of models, with a capability threshold below which it disappears. & Both & Estimation \\
H4 & There is a time scale at which the agentic system's thinking time starts to cost more than its coordination gains. & Both & Estimation \\
\bottomrule
\end{tabularx}
\end{table}

H1 alone is not a positive result, and H2 alone does not support the
registered claim either. A system that wins on novel disruptions but loses
clearly on routine ones has not matched the baseline first. Both have to hold.

The pre-registration also says in advance what would count as failure. H1
fails if the paired mean survival difference is at or below $-\delta_{NI}$. H2
fails if there is no stochastic dominance on the novel set, or if the median
level difference is below $\delta_m$. A third outcome is named as well, the
\emph{null zone}: both systems collapse to level 0 or 1 across the novel set.
That outcome is reported as an uninformative campaign, not as a tie.

One boundary applies to every answer this design gives. It tests whether the
agentic system won \emph{on these scenarios}. It cannot show whether such a
system can ever win. Whether the scenarios actually demanded a change of plan
is a question the design assumed rather than tested, and
Chapter~\ref{ch:critique} returns to it.

% =====================================================================
\section{Parity Between the Two Conditions}
\label{sec:eval:parity}

A comparison is only fair if both systems see the same world and can say the
same things to it. Parity is audited on five axes: the information horizon,
the action vocabulary, the world model, the scenario inputs, and the output
substrate. It is demonstrated by a capability table and an executable
conformance test, not asserted. Every known deviation is declared with its
direction, including two that favour the control. Chapter~\ref{ch:validity}
carries the full audit and does not need repeating here.

One difference is not a deviation and is named here so that it is not mistaken
for one. Every scenario declares bus-2 as a reserve. The baseline sends it out
in parallel from its first schedule. The agentic side holds it back unless one
of four named reasons releases it. Both sides have the same capability, so this
is a difference in \emph{decisions}, which is exactly what the comparison
measures. Section~\ref{sec:val:reserve} describes the mechanism.

% =====================================================================
\section{Scenario Suite}
\label{sec:eval:scenarios}

The suite has two arms (Table~\ref{tab:eval:scenarios}). The content of each
disruption is described in Section~\ref{sec:cs:disruptions}. Every scenario
exists as a pair of files, one per condition, read by the same loader under
the same hazard.

\begin{table}[htbp]
\centering
\caption{The scenario suite as evaluated.}
\label{tab:eval:scenarios}
\small
\begin{tabularx}{\textwidth}{ll>{\raggedright\arraybackslash}X}
\toprule
Arm & ID & Scenario \\
\midrule
Anticipated & S01 & \code{closed_loop}, the undisrupted control \\
            & S02 & \code{road_closure}, a road closes during transport \\
\midrule
Novel & S04 & \code{medical_emergency_in_flight} \\
      & S05 & \code{unverified_persons_report} \\
      & S06 & \code{conflicting_evacuation_orders} \\
      & S07 & \code{bus1_suspension_capacity_reduction} \\
      & S08 & \code{sc1_partial_capacity_loss} \\
\midrule
Gate only & (S03) & \code{baseline_shelter_closure}, a fixture for the negative control. Never scored. \\
\bottomrule
\end{tabularx}
\end{table}

\paragraph{The anticipated arm.} These are disruptions the baseline has a
written branch for, and both were visible during development. A third
scenario, a shelter closure, had no agentic counterpart. A scenario cannot be
reported for one side only, so it was excluded from scoring on 2026-09-10. It
lives on as the fixture for the negative control
(Section~\ref{sec:eval:noreplan}). The reason it has no agentic pair is a
property of the implementation. The live shelter agent recomputes its free
beds from its own identity on every query and has no handler that lowers them.
A capacity disruption would therefore invert: a coordinator that correctly
redistributed would score below one that ignored the loss. This leaves an
anticipated arm of two scenarios, which is thin for a per-scenario analysis.
It is declared as a limitation.

\paragraph{The novel arm.} Five disruptions were supplied blind and revealed
only after the baseline was frozen (Section~\ref{sec:eval:firewall}). Their
identifiers were reserved as \code{novel_1} to \code{novel_5} before their
content existed, so the set could not be extended quietly after results were
seen. Every supplied scenario had to be run. A novel scenario may not be
dropped for being too hard, badly formed, or unfavourable.

\paragraph{What is not in the suite.} An earlier plan had a third tier, a
scenario-scaling sweep over resident count, shelter topology and fleet size.
It was registered as exploratory and was never run. The evidence therefore
rests on one case study, and that is stated as a limitation rather than
argued around.

% =====================================================================
\section{The Disruption Firewall}
\label{sec:eval:firewall}

The novel disruptions are the part of the design an author could most easily
bias without meaning to. The firewall exists to prevent that. The disruptions
are written by the external supervisor, who is a domain expert, and they are
written blind to the baseline. They are revealed only after the baseline is
committed and tagged in git. The tag and its timestamp are the auditable proof
that the baseline could not have encoded them in advance.

The supervisor received a \emph{characteristics-only brief}. It named the kind
of disruption wanted (compound, involving several actors, unfolding in the
middle of the mission, not a single clean road closure) and never the content.
The brief is reproduced in Appendix~\ref{app:brief}. Its purpose is to steer
blind content towards disruptions that are hard enough to matter, without the
author seeing any of them.

\paragraph{As executed.} No separate baseline-freeze tag was cut. The freeze is
commit \texttt{67a14f8}, made at 11:41 on 2026-09-10 and tagged
\code{eval-day-2026-09-10}. The five novel scenarios were supplied through the supervisor and
saved the same afternoon, between 15:27 and 16:04, and then committed. The
freeze therefore precedes the reveal by several hours. One later change to
baseline code was merged after the reveal (\texttt{e9e20cd}). It added the
switch that turns off re-planning for the negative control. Amendment~6
discloses it, together with the evidence that the doctrine's own behaviour did
not change.

\paragraph{The residual risk.} Blind supply has a known cost, the null zone. A
blind disruption may land too easy, so both systems tie, or too hard, so both
collapse, and there is no time to supply a new one. The brief was the only
mitigation. A behavioural pilot on self-written held-out disruptions was
considered and declined. This risk was recorded before the evaluation, not
explained afterwards.

% =====================================================================
\section{Primary Outcome}
\label{sec:eval:primary}

The primary outcome is the Recovery-Outcome Level, an ordinal scale from 0 to 4.
It is built from a continuous quantity, the survival rate $S$, which is the
co-primary outcome for safety. Both are computed by a frozen function of the
event log. This section builds them up in order: the denominator, the
numerator, and then the ladder.

\subsection{Feasibility Ceiling}
\label{sec:eval:ceiling}

Not every disruption leaves every resident savable. If a shelter loses beds,
the best possible result is lower. The outcome is therefore normalised by the
number of residents who \emph{could} have been saved, the feasibility ceiling
\begin{equation}
N =
\begin{cases}
\min\!\left(P,\ \sum_s b_s\right) & \text{if } \sum_s b_s > 0,\\[2pt]
P & \text{if no shelter declared any capacity,}
\end{cases}
\label{eq:eval:ceiling}
\end{equation}
where $P$ is the population the scenario declares and $b_s$ is the number of
beds shelter $s$ declared. The population cap is there because evacuating more
residents than exist is not a higher achievement. The bed cap is there because
residents cannot be sheltered in beds that do not exist. The fallback is not
``zero beds''. With no declaration anywhere the sum is unknown, and treating it
as zero would make every run trivially complete. The fallback is written into
the run's notes and is never silent.

The ceiling is a known point of attack. If it could depend on what a run did,
a system could lower its own denominator and score better by doing worse. The
rule is therefore that the ceiling comes from what the scenario and the
shelters \emph{declared}, never from the run's own placements.

The idea of normalising an ordinal outcome by what is achievable has a
precedent in clinical trials. In stroke and brain-injury trials, a minority of
studies define a favourable outcome relative to how severe the case was at the
start, because a fixed threshold is out of reach for the most severe patients
\parencite{murray2005,nunn2016}.
That literature also names the cost, and it applies here too. A normalised
outcome is harder to interpret and harder to compare across studies.

\subsection{Survival Rate and Loss Decomposition}
\label{sec:eval:survival}

Residents in the simulation are not killed. They are lost to the operation, and
there are exactly three ways for that to happen. A resident is \emph{never
moved}: still in the building when the zone floods, after which nobody can
board. A resident is \emph{late}: set down after the flood arrival of the zone,
which means in the water and not in a shelter. Or a resident is \emph{refused}:
delivered to a shelter that was closed or full. The third channel is the
``routed into danger'' case, and it is why safety cannot be read from timing
alone. A system can be perfectly on time and still deliver 60 residents to a
shelter with no beds.

Two filters decide who counts, and they are applied in this order. First the
deadline: an arrival counts only if it was set down at or before the flood
arrival of its zone. Then an admission pass: the in-time arrivals are taken in
time order, and each shelter admits residents only up to its declared beds.
Only the excess over the beds is invalidated, never the whole delivery and
never the whole run. This rule exists so that ``saved everyone, with one flaw''
cannot rank below ``cleanly saved a third''.

The residents who pass both filters are the validly sheltered residents $V$.
The survival rate is
\begin{equation}
S = \frac{V}{N},
\label{eq:eval:survival}
\end{equation}
the share of the residents who could have been saved that were set down, in
time, at an open shelter with a free bed. Every run also reports where the
others went:
\begin{equation}
N = \underbrace{(N - H)}_{\text{never moved}} + \underbrace{M}_{\text{late}} + \underbrace{O}_{\text{refused}} + \underbrace{V}_{\text{safe}},
\label{eq:eval:loss}
\end{equation}
where $H$ is the number of residents reported sheltered, $M$ the number set
down late, and $O$ the number refused for want of a bed. $S$ alone cannot tell
``too slow'' apart from ``drove them somewhere lethal''. These are different
failures of a coordinator, so the decomposition is always published and never
folded into one number.

\subsection{Recovery-Outcome Level}
\label{sec:eval:rubric}

The primary dependent variable places each run on a five-level ladder. It has
one axis, valid completeness $C = S$, and it uses bus-kilometres $K$ only to
split the top two levels:
\begin{equation}
L(C, K) =
\begin{cases}
0 & C < 0.10 \quad \text{collapse}\\
1 & 0.10 \leq C < 0.50 \quad \text{major loss}\\
2 & 0.50 \leq C < 1.00 \quad \text{partial}\\
3 & C = 1.00 \ \wedge\ K > 1.25\,K^\ast \quad \text{complete, degraded}\\
4 & C = 1.00 \ \wedge\ K \leq 1.25\,K^\ast \quad \text{complete and efficient}
\end{cases}
\label{eq:eval:ladder}
\end{equation}
Figure~\ref{fig:eval:ladder} draws it.

% The Recovery-Outcome Level on the completeness axis, with the sensitivity
% cuts. Source: EVALUATION_PROTOCOL sec.6 and sec.6.3, thresholds.py.
% Input from inc/evaluation.tex, Section "Recovery-Outcome Level".
\begin{figure}[htbp]
\centering
\begin{tikzpicture}[x=11cm, y=1cm,
  band/.style={draw, thick, minimum height=1cm, align=center, font=\footnotesize},
  lab/.style={font=\footnotesize},
]
  % Bands on the completeness axis C in [0, 1).
  \draw[band, fill=red!18]    (0.00,0) rectangle (0.10,1);
  \draw[band, fill=orange!18] (0.10,0) rectangle (0.50,1);
  \draw[band, fill=yellow!22] (0.50,0) rectangle (1.00,1);
  \node[lab, align=center] at (0.05,0.5) {\textbf{0}};
  \node[lab, align=center] at (0.30,0.5) {\textbf{1} \ major loss};
  \node[lab, align=center] at (0.75,0.5) {\textbf{2} \ partial};
  \node[lab, above] at (0.05,1.0) {collapse};

  % The point C = 1.00, split by bus-km against the clean-run reference.
  \draw[band, fill=green!12] (1.04,0.52) rectangle (1.30,1.12);
  \draw[band, fill=green!30] (1.04,-0.12) rectangle (1.30,0.48);
  \node[lab, align=center] at (1.17,0.82) {\textbf{3} \ degraded\\[-2pt] $K > 1.25\,K^\ast$};
  \node[lab, align=center] at (1.17,0.18) {\textbf{4} \ efficient\\[-2pt] $K \leq 1.25\,K^\ast$};
  \draw[-{Stealth[length=2mm]}, thick] (1.00,0.5) -- (1.04,0.5);
  \node[lab, above] at (1.17,1.12) {$C = 1.00$};

  % Axis and ticks.
  \draw[-{Stealth[length=2mm]}] (0,-0.35) -- (1.02,-0.35) node[right, lab] {$C$};
  \foreach \x/\t in {0/0, 0.10/0.10, 0.50/0.50, 1.00/1.00}
    \draw (\x,-0.28) -- (\x,-0.42) node[below, lab] {\t};

  % Sensitivity cuts for the level-1/2 boundary.
  \foreach \x in {0.40, 0.60}
    \draw[dashed, gray!70!black] (\x,-0.1) -- (\x,1.25);
  \node[lab, text=gray!40!black, align=center] at (0.50,1.55)
    {sensitivity cuts 0.40 / 0.50 / 0.60};
\end{tikzpicture}
\caption[The Recovery-Outcome Level]{The Recovery-Outcome Level $L$ as a
function of valid completeness $C$. Levels 0 to 2 are read off $C$ alone.
Only a complete run ($C = 1.00$) is split, by its bus-kilometres $K$ against
the clean-run reference $K^\ast$. The dashed lines mark the alternative
level-1/2 cuts used in the sensitivity check.}
\label{fig:eval:ladder}
\end{figure}


\paragraph{Why each cut sits where it does.} Each cut has a domain reason, and
each reason was written down before any result existed. The cut at 0.50 is
``half the home''. An evacuation that gets fewer than half the residents out
has failed in a different way from one that got most of them out and fell
short. The first is a mass-casualty event. The second is an incomplete but
largely successful operation. The cut at 0.10 is one part-load of a single
50-seat bus in a 100-resident home. A run that never moves more than that did
not perform an evacuation, it failed to start one. The 25\,\% band at the top
is round headroom. Handling a disruption legitimately costs extra driving,
such as a detour or a second pickup, and a tighter band would score necessary
work as waste.

\paragraph{The clean-run reference.} $K^\ast$ is the median bus-kilometres of
the control runs in the same cell that delivered everyone with per-shelter
attribution. It is defined only when at least three such runs exist. Only
complete runs qualify, because a run that sheltered part of the home also drove
less, and admitting it would tighten the band for everyone else. Three is the
smallest number of runs for which a median is not just one observation. When
$K^\ast$ is undefined, no run may rise above level 3, and the reason is written
into its notes. The comparison is always within a system: it asks whether a
system drove further than it needed to, never whether it drove further than
the other system.

\paragraph{Where safety and efficiency meet.} Level 4 is the one place where the
ladder mixes safety with efficiency. This is exactly the kind of blend the
gate structure of Section~\ref{sec:eval:gates} warns against. The protocol
therefore states it openly as an identity,
\begin{equation}
L = 4 \iff \text{Gate 1 passed} \wedge \text{Gate 2 passed},
\end{equation}
so that it is visible rather than hidden inside the ladder.

\paragraph{Why an ordinal ladder.} A pass/fail outcome would throw away the
difference between losing one resident and losing fifty. The design borrows
from stroke trials, where the modified Rankin Scale is an ordered outcome with
a long methods literature. That literature recommends keeping the ordinal
structure instead of cutting it into two
\parencite{optimisinganalysisofstroketrialsoastcollaboration2007a,ganesh2018},
and it documents disagreement about where a single cut should fall
\parencite{nunn2016}. The sensitivity check in Section~\ref{sec:eval:sensitivity}
answers that second point. The analogy is \emph{structural} only. The ladder
borrows the shape of an ordered clinical scale and its analysis methods. It
makes no claim that an evacuation outcome is clinically comparable to a
stroke outcome. The same literature raises one counter-point, which is
accepted here: an ordinal analysis treats every step as equal in rank, so a
step from 0 to 1 counts the same as a step from 3 to 4. The analysis never
does arithmetic on levels (Section~\ref{sec:eval:analysis}), so it never has to
say how far apart two steps are.

\paragraph{Why an outcome and not an area.} Resilience engineering usually
measures ``graceful degradation'' as the area lost under a functionality curve,
the resilience triangle \parencite{bruneau2003}. The ladder deliberately does
not. An area weights every moment equally, so very different trajectories can
receive the same score \parencite{cremen2025}.
In an evacuation, what matters is whether residents were in a shelter before
the water arrived. The ladder scores that outcome directly.

\subsection{Deterministic Scorer}
\label{sec:eval:scorer}

No human scores anything. A frozen function of the event log computes every
outcome, and the same function runs on both systems. It reads the log
described in Section~\ref{sec:data:evalbackbone}. This is a deliberate contrast
with the common practice of using a language model as a judge
\parencite{zheng2023}, and with agent benchmarks that use one for some rubric
items \parencite{xu2025}. A model judging models would bring back exactly the
nondeterminism the design is trying to measure.

Two conventions run through the scorer. First, a missing value is not a zero.
\texttt{None} means ``this run's log cannot answer that question'' and is never
treated as 0. The worst failure of a thesis instrument is not a wrong number,
it is a number that looks like a measurement when nothing was measured.
Second, every interpretive constant lives in one file, \code{thresholds.py},
which contains data and no logic. Its hash is written into every run's
manifest. The claim ``the thresholds were frozen before the runs'' is
therefore a checkable fact about one file's history, not a promise about a
whole package.

% =====================================================================
\section{The Three Gates}
\label{sec:eval:gates}

A system is not better because it wins a combined score. A combined score lets
a system trade safety for speed. The protocol therefore reads three gates in a
fixed order, and a gate is read only if every gate above it passed
(Figure~\ref{fig:eval:gates}). The order is the point. A clearance time
measured over a run that abandoned part of the home is not a slow number, it is
an undefined one. Reading speed without safety is what lets a system look fast
by giving up on the hardest residents. The gate order makes that impossible by
construction.

% The three hierarchical gates, where each hypothesis is read, and where the
% no-replan control sits. Source: EVALUATION_PROTOCOL sec.0, sec.3-5, sec.5.4.
% Input from inc/evaluation.tex, Section "The Three Gates".
\begin{figure}[htbp]
\centering
\begin{tikzpicture}[
  gate/.style={draw, thick, rounded corners=2pt, fill=blue!6,
    minimum width=5.0cm, minimum height=1.25cm, align=center, font=\small},
  side/.style={draw, rounded corners=2pt, fill=gray!8, align=left,
    font=\footnotesize, text width=3.6cm, inner sep=4pt},
  ctrl/.style={draw, dashed, rounded corners=2pt, fill=orange!8, align=left,
    font=\footnotesize, text width=3.6cm, inner sep=4pt},
  fail/.style={font=\footnotesize, text=red!60!black, align=left},
  flow/.style={-{Stealth[length=2.2mm]}, thick},
  link/.style={-{Stealth[length=1.8mm]}, gray!70!black},
]
  \node[gate] (g1) {\textbf{Gate 1: Safety}\\ Did the residents survive? \quad $S = V/N$};
  \node[gate, below=1.1cm of g1] (g2) {\textbf{Gate 2: Speed and Flow}\\ How quickly? \quad $T_{50}, T_{90}, T_{95}, T_{100}$};
  \node[gate, below=1.1cm of g2] (g3) {\textbf{Gate 3: Adaptability}\\ Does an advantage hold\\ under novel shocks?};

  \draw[flow] (g1) -- node[right, font=\footnotesize] {passed} (g2);
  \draw[flow] (g2) -- node[right, font=\footnotesize] {passed} (g3);

  \node[fail, left=0.35cm of g1.south west, anchor=north east, yshift=-0.05cm]
    (f1) {fail: run is a failure,\\ Gate 2 and 3 not read};
  \draw[link, red!60!black] ([xshift=-0.6cm]g1.south) |- (f1.east);

  \node[side, right=0.6cm of g1] (h1) {\textbf{Anticipated arm} (H1)\\ non-inferiority:\\ $RD > -\delta_{NI}$, \ $\delta_{NI}=0.05$};
  \node[side, right=0.6cm of g3] (h2) {\textbf{Novel arm} (H2)\\ superiority on $S$, then the\\ effect rule: median shift\\ $\geq \delta_m = 1$ level \emph{and} dominance};
  \draw[link] (h1.west) -- (g1.east);
  \draw[link] (h2.west) -- (g3.east);
  \draw[link] (h2.north west) -- ([yshift=-0.3cm]g1.south east);

  \node[ctrl, left=0.6cm of g3] (nr) {\textbf{No-replan control}\\ campaign gate: must score\\ strictly below the baseline\\ on $\geq 1$ disruption};
  \draw[link, dashed] (nr.east) -- (g3.west);

  \node[font=\footnotesize, text=gray!50!black, below=0.25cm of g2.south east, anchor=north west, xshift=0.15cm]
    {$L = 4 \equiv$ Gate 1 $\wedge$ Gate 2};
\end{tikzpicture}
\caption[The three gates of the evaluation protocol]{The three gates of the
evaluation protocol. A gate is read only if every gate above it passed.
Hypothesis H1 is read at Gate~1 on the anticipated arm. Hypothesis H2 is read
on the novel arm, first as superiority on survival and then as a minimum effect
on the level scale. The no-replan control is a negative control on the whole
campaign, not a third system under test.}
\label{fig:eval:gates}
\end{figure}


The reference protocol this design descends from has four gates: safety,
operational efficiency, equity, and robustness. This thesis runs three. Equity
is not dropped. It is registered as undefined, and the reason is printed in the
same place as the other gates' verdicts (Section~\ref{sec:eval:metrics}).

\subsection{Gate 1: Safety}
\label{sec:eval:gate1}

Gate 1 asks whether the residents survived. Its effect measure is the risk
difference, computed paired, per scenario, treatment minus control:
\begin{equation}
\widehat{RD} = \frac{1}{n}\sum_{i=1}^{n} \big(S_{A,i} - S_{B,i}\big),
\label{eq:eval:rd}
\end{equation}
over the $n$ matched pairs. A positive value favours the agentic system. It is
always reported with a bootstrap confidence interval and a Wilcoxon check
(Section~\ref{sec:eval:analysis}). The pass rule depends on the arm:

\begin{itemize}
  \item \textbf{Anticipated arm:} pass when $\widehat{RD} > -\delta_{NI}$, with
    $\delta_{NI} = 0.05$. The agentic system is not worse on survival by more
    than the margin.
  \item \textbf{Novel arm:} pass when $\widehat{RD} > 0$ \emph{and} the paired
    bootstrap interval excludes zero. The agentic system is better on
    survival.
\end{itemize}

If no margin has been registered, the verdict is \texttt{unregistered}, and
that is not a pass. A run or condition that fails Gate 1 is reported as a
failure. Its Gate 2 and Gate 3 figures are not published as achievements.

\paragraph{The two margins.} $\delta_{NI} = 0.05$ is five residents out of a
home of 100. It is stated on $S$ because a non-inferiority margin must be a
quantity of harm: ``how many residents may we lose and still call it
matching?'' has an answer, while the same question in levels does not. Five is
half of the smallest unit the ladder already recognises, the 0.10 part-load.
The margin is strict on purpose. A wide margin makes ``matching'' easy to show
and worthless to claim. $\delta_m$, the minimum effect for the novel arm, is
one whole level, because on an ordinal scale the level \emph{is} the
substantive unit. $RD$ lives between $-1$ and $1$ while $\delta_m$ is one level,
so the two are never read against each other. Reading them together would make
every possible novel win look too small. The code keeps them apart, and a test
enforces it.

Passing Gate 1 on the novel arm is not yet a claim to have won. Whether a win
is large enough to be worth claiming is a separate registered rule, stated on
the level scale (Section~\ref{sec:eval:analysis}).

\subsection{Gate 2: Speed and Flow}
\label{sec:eval:gate2}

Gate 2 asks how quickly the residents reached safety, given that they
survived. It is read over the Gate-1-clearing runs only. Its main measures are
clearance times at four quantiles, $q \in \{0.50, 0.90, 0.95, 1.00\}$. With
the admitted arrivals in time order and $F(j)$ the running count of admitted
residents,
\begin{equation}
T_q = \min\big\{\, t_j \ :\ F(j) \geq \lceil q\,N \rceil \,\big\}.
\label{eq:eval:tq}
\end{equation}

The single most important choice in Gate 2 is the denominator. The quantile
$q$ is a share of the ceiling $N$, not of the residents a run happened to save.
If a run's survival is below $q$, then $T_q$ does not exist. It is reported as
censored. It is never imputed, never dropped silently, and never replaced by
the time of the last resident who did arrive. A run that abandons 15\,\% of
the home has no $T_{90}$ and no $T_{95}$, and it cannot appear in a
clearance table as a fast run. Defining $q$ over survivors would bring back
exactly the gaming the quantiles exist to prevent: a system that saves the
easy half and stops would post an excellent $T_{90}$ of that half.

``Flow'' in the reference protocol means road occupancy and queues. This
simulator has no congestion model, so flow is measured as operational effort
instead: bus-kilometres, residents per bus-kilometre, and the response-time
cost of Section~\ref{sec:eval:responsetime}.

Gate 2 is descriptive, not pass/fail. One expectation was registered in
advance: the baseline would win on latency and cost. That was stated as the
honest trade-off before the runs, not discovered afterwards.

\subsection{Gate 3: Adaptability}
\label{sec:eval:gate3}

Gate 3 asks whether an advantage survives disruptions that nobody planned for.
It has three sub-tests.

\begin{enumerate}
  \item[(a)] \textbf{The anticipated precondition.} The baseline must pass
    100\,\% of the anticipated set. A control that cannot run a clean
    evacuation makes every comparison meaningless. On this set the agentic
    system is tested only for non-inferiority.
  \item[(b)] \textbf{The novel set.} The per-scenario level distributions over
    the $k$ replicates are compared by first-order stochastic dominance. This
    is the primary test of H2.
  \item[(c)] \textbf{The degradation curve.} The advantage
    $A(z) = Y_A(z) - Y_B(z)$, where $Y$ is the \emph{median} level at time
    scale $z$, and the crossover band where the advantage first vanishes.
\end{enumerate}

Sub-test (c) uses the time scale \code{SIM_TIME_SCALE} as its axis. It asks at
what exchange rate between wall-clock and simulated time the agentic system's
thinking time starts to cost residents. It needs runs at several scales. Only
one scale, 8.4, was executed, so the curve has one point per cell and the
crossover was never read. H4 was therefore not tested
(Section~\ref{sec:res:h4}).

\subsection{The No-Replan Control}
\label{sec:eval:noreplan}

A third condition runs beside the two under test. It is the doctrine baseline
with its re-planning switched off. The only path to a second dispatch order is
disabled, so the first committed plan is its only plan. Everything else in the
doctrine stays unchanged. It is the one ablation this design carries, and it
serves two purposes. It separates \emph{re-planning} from \emph{reasoning about
what to re-plan}. And it is the instrument's negative control, under a rule
registered as a gate on the whole campaign:

\begin{quote}
Before any result is claimed, the no-replan control must score strictly below
the doctrine baseline on at least one disruption.
\end{quote}

If a system that ignores every disruption scores the same as one that handles
them, the instrument is not measuring adaptation, and no comparison built on it
means anything.

\paragraph{Where the gate was read.} The gate was read on the shelter-closure
fixture, a scenario outside the scored set. On that fixture a re-plan changes
who is delivered. It passed on 2026-09-14: the doctrine reached level 3 with
$S = 1.0$, and the no-replan control reached level 0 with $S = 0.0$. The
protocol also records, in writing, why the gate was not read on the scored
scenarios. On those files the control ``cannot differ from the doctrine in
outcome, by construction''. It commits the same orders on the undisrupted run
and on the novel files, and on \code{road_closure} the completeness ties
whatever was ordered. That note was written as the reason to move the gate to a
fixture. Its consequence for the scored arms was not drawn at the time.
Chapter~\ref{ch:critique} draws it (Section~\ref{sec:crit:gate0}).

% =====================================================================
\section{Secondary Instruments}
\label{sec:eval:secondary}

\subsection{Held-Out Decision Rubric}
\label{sec:eval:heldout}

Four of the five novel scenarios say in their own text that their main effect
is ``judged mainly on the decision rubric''. The simulator models no patient
(S04), no staircase or police officer (S06), no payload cap (S07), and no
partly damaged shelter (S08). A coordinator that ignores the report can still
reach a full level. S05 is not flagged that way, but it carries an explicit
checklist and is judged beside the others. A decision rubric was therefore
added. It was decided on 2026-09-11, with anticipated-arm pilot data in hand and
no novel-arm data, and registered by Amendment~6 as the novel arm's
\emph{secondary} outcome.

The rubric is secondary and descriptive. It is never confirmatory, never turned
into a level, and never merged with $L$. It is reported as counts of yes, no
and unanswerable per criterion, by disruption, system and model. It reads
\emph{structure only}: who was addressed and when, what the orders booked, and
what was set down where. It never reads message text or model reasoning. A
clause that depends on what a message \emph{says} is listed as not judged
rather than approximated with keywords. An unanswerable criterion stays
unanswerable and is never counted as a failure.
Table~\ref{tab:eval:heldout} condenses the 16 criteria.

\begin{table}[htbp]
\centering
\caption{The held-out decision rubric, condensed. Each clause from the
supplier's scenario text becomes a mechanical rule on the log.}
\label{tab:eval:heldout}
\small
\begin{tabularx}{\textwidth}{l>{\raggedright\arraybackslash}X}
\toprule
Scenario & Criteria \\
\midrule
S04 & A medical team is requested to meet bus-1 at its destination shelter. Bus-1's leg is neither withdrawn nor diverted. ZKD is informed. \\
S05 & The report is verified with police or ZKD before it is retracted. The standing runs land as ordered. The reserve bus is committed as a hedge by 10:57. Nothing new is still moving after the retraction. \\
S06 & The nursing home is told by 10:58. The conflict is escalated to ZKD. Seats on the standing legs still cover the whole population. \\
S07 & No bus-1 leg sets down more than 25 after the report. Bus-2 is committed. Seats still cover the population. \\
S08 & No order books SC1 above 40. The excess goes to SC2. The valid count, with SC1 capped at 40, reaches the population. \\
\bottomrule
\end{tabularx}
\end{table}

What the baseline would do was written down before any held-out run. The
doctrine answers every held-out report with the same hold: status messages to
ZKD, the transport operator and the home, and no new order. It therefore
passes the ``was informed'' criteria by construction and fails the ones that
need a changed plan. That is a finding about the doctrine, not a flaw in the
rubric. It is also part of why the rubric stays descriptive. It shows
\emph{where} the two conditions differ, not an effect size H2 could read.

\subsection{Secondary and Diagnostic Metrics}
\label{sec:eval:metrics}

Beside the primary outcome, each run reports the secondary measures in
Table~\ref{tab:eval:kpis}. None of them can move the level.

\begin{table}[htbp]
\centering
\caption{Secondary measures, in the registered order of priority after the
clearance quantiles and the loss decomposition.}
\label{tab:eval:kpis}
\small
\begin{tabularx}{\textwidth}{l>{\raggedright\arraybackslash}X}
\toprule
Measure & Meaning \\
\midrule
Bus-km $K$ & Distance actually driven. Charged once per leg, priced from where the bus stands. \\
Resource efficiency & $\eta = V / K$, validly sheltered residents per bus-km. \\
Wall latency, tokens & The response-time cost (Section~\ref{sec:eval:responsetime}). \\
Capacity overloads & Distinct events where a shelter received more than it declared. These do affect the outcome, through $O$. \\
Constraint violations & Rejected commands and rejected messages, counted in two separate channels. \\
Time to safety & Time of the last counted dropoff. This is \emph{not} $T_{100}$, which is censored unless everyone arrived. \\
\bottomrule
\end{tabularx}
\end{table}

Three points about these measures matter for reading the results.

\paragraph{Bus-km is not comparable across systems.} Each total is consistent
within its own system, but they are composed differently. The agentic total
sums the legs a driver committed to. The baseline total is priced from depot to
pickup to shelter for each assignment, because its commander allocates rather
than drives. That is enough for the level-3-versus-4 split, which is a
within-system comparison. It is not enough for a cross-system efficiency
claim, and the code refuses to present it as one.

\paragraph{Two kinds of invalid action.} A \emph{command} violation is a verb
the world has no handler for. An \emph{outbound} violation is a message to a
recipient outside the address book, or of a type the agent may not send. The
second is the coordination violation, and it bears on the shared-database
architecture. Each rejection carries a reason code from a closed list, so the
counts survive a change in the wording of an error message.

\paragraph{Diagnostics, not scores.} Leg count and supersede count explain an
outcome, for example ``many plan changes and a low level'' reads as thrashing.
They are never used to score success, because a supersede count rewards
churn.

\paragraph{Equity is registered as undefined.} Humanitarian logistics usually
treats equity as an outcome on the same footing as efficiency and
effectiveness \parencite{huang2012}.
An equity gate would compare survival across subgroups. Here it is registered
as undefined, and the reason is a property of the case. A nursing home is a
homogeneous vulnerable population. The whole population \emph{is} the
vulnerable group, residents are counted as integers with no attribute to split
on, and a difference across one group is always zero. A better measurement was
considered: mobility tiers such as walking, wheelchair and bed-bound, with an
equity margin. It was declined because it was new scope competing with the
runs, and because a subgroup structure introduced late would be one chosen with
the campaign in view. The ladder therefore cannot tell a run that sheltered the
60 most mobile residents from one that sheltered the 60 least mobile. This is
declared as a limitation.

\subsection{Response-Time Cost}
\label{sec:eval:responsetime}

The agentic system thinks, and thinking takes time. That cost is reported as a
trade-off the baseline was expected to win, not as an axis the agentic system
claims. It has two views. The first is the raw wall-clock latency $\lambda$ of
the model calls, per model. The second is the operational delay in simulated
time, $\Delta_{\text{op}} = \lambda \cdot \sigma$, where $\sigma$ is the time
scale. It is reported as a fraction of the disruption's own decision budget,
$\phi = \Delta_{\text{op}} / (D_z - t_{\text{first event}})$. The same delay can
be harmless under a slow disruption and fatal under a tight deadline, and the
fraction makes that visible.

Two cautions apply. $\lambda$ is summed over every turn of every agent, but the
ten agents think in parallel, so $\Delta_{\text{op}}$ is an upper bound on the
delay the operation actually suffered. And both quantities are withheld unless
$\sigma$ is supplied explicitly as a logged constant, so no delay is ever
computed against an assumed scale.

The primary outcome deliberately leaves latency and cost out. The humanitarian
logistics literature has an established argument that operational cost is not
the outcome that matters in a disaster \parencite{holguin-veras2013}.

% =====================================================================
\section{Statistical Analysis}
\label{sec:eval:analysis}

\paragraph{No arithmetic on levels.} The levels are ranked, not spaced. The
distance from 0 to 1 is not the distance from 3 to 4, so a mean level is not a
quantity that exists. Levels are compared by distributions and medians only.
Where a continuous quantity is needed, $S$ is used, since the level is a
function of it. $RD$ is a mean of paired differences precisely because $S$ is
continuous.

\paragraph{Stochastic dominance.} Level distributions are compared by
first-order stochastic dominance, which needs no arithmetic on an ordinal
scale \parencite{yalonetzky2013,davidson2000,davidson2013}. With
$\hat p_{\mathcal A}(L)$ the share of runs in sample $\mathcal A$ that score
above level $L$,
\begin{equation}
\mathcal A \succ \mathcal B \iff
\big(\exists L : \hat p_{\mathcal A}(L) > \hat p_{\mathcal B}(L)\big)
\wedge
\big(\nexists L : \hat p_{\mathcal B}(L) > \hat p_{\mathcal A}(L)\big).
\label{eq:eval:dominance}
\end{equation}
The verdict is one of: $\mathcal A$ dominates, $\mathcal B$ dominates, neither
(the distributions cross), identical, or undetermined (an empty sample). The
single-ladder comparison used here is the one-dimensional special case of the
ordinal dominance conditions in \textcite{yalonetzky2013}. Dominance checks
lose power in sparse tails, and with $k = 10$ and five levels some cells are
nearly empty. The raw level counts are therefore always printed beside the
verdict.

\paragraph{Confidence intervals and the paired check.} Every risk difference
carries a percentile bootstrap interval. The paired differences are resampled
with replacement 10\,000 times at a fixed seed (20260904), and the 2.5th and
97.5th percentiles of the resampled means form the 95\,\% interval. An empty
sample gets no interval rather than a meaningless $(0,0)$. A Wilcoxon
signed-rank test on the same differences is reported beside it as an
assumption-light check. Zero differences are dropped, as Wilcoxon's rule
requires, but they are \emph{counted}, because a set of pairs that is mostly
exact ties is a result about the systems. The test is exact for up to 20
non-zero differences without ties, and uses a normal approximation with tie
correction otherwise.

\paragraph{The H1 rule.} H1 passes on a scenario when
$\widehat{RD} > -\delta_{NI}$. This is a rule on the point estimate, and it
carries no significance level. The pre-registration first registered a
one-sided test on the lower bound of the interval. The implementation had
always decided on the point estimate, and on 2026-09-12, with pilot data seen,
the author registered the implemented rule rather than change the code. The
consequence is stated openly. An H1 pass has no controlled error rate, so a
system that sits exactly at the margin passes about half the time at any $k$.
A pass is therefore reported as ``the point estimate lies inside the margin'',
never as ``non-inferiority demonstrated''. The margin itself did not change.

\paragraph{The H2 effect rule.} Gate 1 decides whether the agentic system is
better on survival. A second rule decides whether the win is large enough to
claim, on the level scale. It has two conditions, and both are required:
\begin{equation}
\operatorname{median}\big\{L_{A,i} - L_{B,i}\big\} \geq \delta_m = 1
\quad\wedge\quad
\mathcal L_A \succ \mathcal L_B.
\label{eq:eval:h2}
\end{equation}
Dominance is a requirement here, not a second number beside the median. Four
runs at level 4 and two at level 0 give a median shift of two whole levels
from a system that abandoned the home in a third of its runs. Under this rule
that is a failure, not a double success.

\paragraph{The registered H2 primary, implemented late.} The pre-registration
also named a primary test for H2: the probability that an agentic run reaches
at least the baseline's level, with an exact Clopper--Pearson interval over the
$k$ runs, combined across scenarios by a sign test and corrected with
Holm--Bonferroni within each arm. None of the three was implemented when the H2
verdict was first written. All three were implemented afterwards and agree with
the verdict (Amendment~15). Because they were coded after the data were seen,
that analysis is exploratory. Section~\ref{sec:crit:power} discusses a second
finding that came with it: at $k = 10$ this registered test could not reach
conventional significance whatever the data.

\paragraph{Multiplicity and aggregation.} H1 and H2 are two claims on two
disjoint scenario sets, each with one primary outcome, so no correction is
applied between them. Within an arm, across scenarios, the correction is
Holm--Bonferroni. Runs are never pooled across scenarios, because scenarios
differ in ceiling, fleet and deadline, and a pooled mean would average
quantities that cannot be compared. Secondary outcomes are not corrected and
are reported as descriptive.

\paragraph{Choosing $k$.} The reference protocol forbids an arbitrary number of
replicates, so $k$ was to be sized on a pilot. The pilot was planned as 24
paired cells on the anticipated set with two models. It ran as 22 paired cells,
unbalanced across models after mid-pilot swaps that the pre-registration
discloses. A resampling tool drew each cell's differences with replacement from
its own pilot differences, 10\,000 times, for $k$ from 3 to 200, and estimated
how often the H1 rule would pass. The target was 0.80. No cell reached it at any
$k$, and the pass probability \emph{fell} as $k$ grew, because every pilot mean
already lay below $-\delta_{NI}$. No value of $k$ could buy H1. Widening the
margin was prohibited by the registration. So $k = 10$ was set on budget and
precision grounds (Amendment~1), and H1 was declared not expected to pass
before the campaign ran. The pilot data were used for sizing only and were
excluded from every confirmatory analysis.

\paragraph{Exclusions and censoring.} Only three reasons can exclude a run: a
simulator or harness crash, a harness timeout, or right-censoring. A run is
censored if it stopped less than 900 simulated seconds after its flood deadline,
\begin{equation}
\text{censored} \iff t_{\text{end}} < D_z + 900\,\text{s},
\end{equation}
because residents still on their way at the deadline never had the chance to
arrive. Scoring such a run would record the harness's impatience as the
system's failure. Fifteen minutes is about one shuttle turnaround. Explicitly
\emph{not} admissible are an unfavourable outcome, an outlying latency, a model
that emits malformed plans (that is a result), or a novel scenario that proves
too hard. Every exclusion is counted and published beside the distribution it
qualifies. The event log is append-only, enforced by a database trigger, so a
run cannot be removed afterwards without the removal being visible.

\subsection{Cut-Point Sensitivity}
\label{sec:eval:sensitivity}

The obvious attack on a threshold-based outcome is that the thresholds were
chosen to produce the result. The defence is to show the result at other
thresholds too. Every score carries its own level at three alternative cuts
for the level-1/2 boundary: 0.40, 0.50 and 0.60. The registered rule is that
the headline may not be reported without this sensitivity table.

This is a small specification curve: the result is shown across the set of
defensible choices instead of along one path
\parencite{simonsohn2020,steegen2016}.
Only the level-1/2 cut is swept. A multiverse should contain only the choices
that are genuinely arbitrary \parencite{delgiudice2021}. The 0.10 cut is
anchored to one part-load of one bus and the 1.00 cut is ``everyone'', so
neither is arbitrary. ``Half the home'' is the most defensible of the three,
but it is also the one a reader is most likely to question, so it is the one
swept.

% =====================================================================
\section{Model Sweep}
\label{sec:eval:models}

The design called for at least three language models across a range of
capability. Each is pinned by digest where possible and named on every run. The
sweep tests whether an advantage is robust across models (H3), and whether
there is a capability threshold below which weak models emit invalid plans.

\paragraph{As executed.} The execution departed from the registered selection
rule, and every departure was recorded as an amendment. The anticipated arm ran
glm-5.2, deepseek-v4-pro and minimax-m3 (Amendment~1). After screening two
further models, the author chose a final set of the three strongest models in
the pool, which overrode the model the registered rule would have picked
(Amendment~5). On the novel arm the third slot then moved three more times,
each time after runs had been seen, and settled on qwen3.5
(Amendments~10 to~12). Three models completed the novel arm at 50 runs each:
glm-5.2, deepseek-v4-pro and qwen3.5. Three others hold partial series:
kimi-k3 (11 runs), gpt-oss (4) and kimi-k2.7-code (3). Partial series are
reported on their own and never pooled with the full arms.

Because the third slot changed after data were seen, the novel arm's model set
and any analysis that depends on it are exploratory. glm-5.2 and
deepseek-v4-pro were fixed before the arm ran and are unaffected. Two further
consequences follow for H3. A set of the strongest models is not expected to
show a capability threshold, and the model below the set (minimax-m3) was not
run on the novel arm, so the threshold cannot be observed there. Amendment~12
also records that levels are not comparable across models, only $S$, because
each model's clean-run reference is its own.

\paragraph{Model drift and nondeterminism.} Hosted \texttt{:cloud} model tags
cannot be frozen to a version. The model metadata is captured on every run, so
drift can be detected, but not prevented. Inference itself is not bit-exact
even with a fixed seed \parencite{yuan2025}, and agentic evaluations are
known to vary from run to run for this reason \parencite{bjarnason2026}. This
is what the replicates are for.

\paragraph{No retry.} Neither the planning agents nor the model client retry a
failed or malformed model call. A malformed response is surfaced as a hard
failure. Adding a retry would give a weak model a second attempt that a strong
model never needed, which is exactly the capability difference the sweep is
meant to measure. A retry would have to be a declared experimental factor,
not a default. The cost of this choice is stated as a limitation: one timed-out
call at a critical step can lose a run.

% =====================================================================
\section{Pre-Registration and Amendments}
\label{sec:eval:prereg}

Every analytic choice was to be fixed before the first evaluation run, so that
no choice could be made after seeing a result. This is the discipline that
stops the ``garden of forking paths'', where choices made after seeing data
bias a result without any intent to cheat, and it prevents forming hypotheses
after the results are known \parencite{gelman2014,kerr1998}.
Pre-registration is an established practice in machine-learning research, with
dedicated workshops at NeurIPS \parencite{bertinetto2021}, and there is
practical guidance for what a pre-registration in a language-technology setting
should contain \parencite{miltenburg2021}.

\paragraph{The amendment rule.} After the pre-registration is tagged, nothing in
it is edited. Changes are appended as dated amendments. Each one states what
changed, why, and whether any evaluation data had been seen at the time. An
amendment made after data were seen is not misconduct. It is a normal event in
a long campaign. But it must be visible, and any analysis it touches is
reported as exploratory rather than confirmatory. Sixteen amendments were
written. Table~\ref{tab:eval:amendments} lists the ones that shape how the
results must be read.

\begin{table}[htbp]
\centering
\caption{The amendments that change how a result is read. All sixteen are in
the pre-registration's amendment log.}
\label{tab:eval:amendments}
\small
\begin{tabularx}{\textwidth}{lll>{\raggedright\arraybackslash}X}
\toprule
\# & Date & Data seen & Change \\
\midrule
1  & 09-12 & Pilot & $k = 10$; campaign models; replicates 1000--1009; H1 declared not expected to pass. \\
5  & 09-14 & Anticipated arm & Final model set chosen by the author over the registered rule. \\
6  & 09-14 & Anticipated arm only & Novel arm fixed before any novel score: five scenarios, freeze and reveal dates, the negative-control verdict, the rubric as secondary, the new thresholds hash. \\
10--12 & 09-16 to 09-18 & Novel arm, partly & The third model slot moves three times, ending at qwen3.5. \\
13 & 09-20 & Whole novel arm & qwen3.5 anticipated \code{closed_loop} cells added to complete its level ladder; closing audit disclosures. \\
14 & 09-21 & Everything & Reporting correction: Gate 1 has 27 cells, not 26. No analytic change. \\
15 & 09-21 & Everything & The registered H2 primary implemented late, hence exploratory; reporting fixes. No confirmatory number moves. \\
16 & 09-22 & Everything & Treatment validity and Protocol v2. Exploratory by construction (Chapter~\ref{ch:critique}). \\
\bottomrule
\end{tabularx}
\end{table}

\paragraph{As executed.} The margins $\delta_{NI}$ and $\delta_m$ and the equity
decision were registered on 2026-09-04, against zero evaluation runs, and
committed to \code{thresholds.py} in the same commit. The tags
\code{prereg-v1} and \code{prereg-amendment-1} were cut together on
2026-09-12, on commit \texttt{7a4d7cc}, the commit the anticipated campaign ran
on. The tag did \emph{not} precede the first evaluation run: the pilot ran
before it, and the pre-registration says so. \code{thresholds.py} changed once
after the reveal, when two fenced sections for the negative control and the
decision rubric were merged in \texttt{e9e20cd}. Neither section holds anything
the ladder reads. The threshold hash recorded with the scored runs,
\texttt{d08c1e35}, is identical across all 310 of them. The frozen artifacts
are listed in Appendix~\ref{app:frozen}.

\paragraph{Campaign gates.} Five conditions had to hold before any result
could be claimed. Table~\ref{tab:eval:campaigngates} gives their status.

\begin{table}[htbp]
\centering
\caption{The campaign gates and their final status.}
\label{tab:eval:campaigngates}
\small
\begin{tabularx}{\textwidth}{l>{\raggedright\arraybackslash}X>{\raggedright\arraybackslash}X}
\toprule
\# & Gate & Status \\
\midrule
1 & The baseline freeze predates the novel reveal. & Met. 11:41 against 15:27, 2026-09-10 (Amendment~6). \\
2 & The baseline passes 100\,\% of the anticipated set. & Evidence present: all 20 campaign baseline runs score $S = 1.0$. Never recorded in an amendment. \\
3 & The no-replan control scores strictly below the baseline on at least one disruption. & Passed on the fixture, level 3 against level 0 (Section~\ref{sec:eval:noreplan}). \\
4 & PARITY deviation 7 is closed or re-argued in writing. & Open when the campaign ran. Closed afterwards by re-argument (Section~\ref{sec:val:dev7}). \\
5 & Everything declared out of scope appears in the limitations. & Section~\ref{sec:eval:scope} and Section~\ref{sec:crit:limitations}. \\
\bottomrule
\end{tabularx}
\end{table}

Gate 4 needs one sentence here. The campaign ran while it was still open, which
the gate did not permit. It was closed afterwards by the written argument the
gate allows: the deviation runs against the baseline, so it cannot explain a
baseline win. Had the agentic side won, the gate would have had to be closed
properly first.

% =====================================================================
\section{Construct Validity}
\label{sec:eval:construct}

The ladder is meant to measure how well a coordinator handles a disruption. That
is a construct that cannot be observed directly, and the ladder is a model of
it. Following \textcite{jacobs2021}, the question of whether the model measures
the construct splits into seven parts.

\paragraph{Face validity.} The ladder speaks the language of the domain:
residents in a shelter before the water arrives. The doctrine baseline it is
compared against was co-certified by the domain expert. There is no empirical
loss data for this facility to anchor it against, and the expert's review is
the substitute.

\paragraph{Content validity.} The ladder covers completeness, the deadline,
shelter capacity and driving effort. It does not cover equity
(Section~\ref{sec:eval:metrics}), the medical care given on the way, or the
content of messages. The decision rubric covers part of the last gap, and the
rest is declared.

\paragraph{Convergent validity.} Measures that should move together, do they?
This is where the diagnostics earn their place. A high supersede count
together with a low level should read as thrashing. The results chapter tests
this link (Section~\ref{sec:res:diagnostics}).

\paragraph{Discriminant validity.} The ladder must separate a coordinator that
handles a disruption from one that ignores it. The negative control exists to
show this, and on the fixture it did. Whether it held on the scored scenarios
is the question Chapter~\ref{ch:critique} answers. It is the weakest part of
this design.

\paragraph{Predictive validity.} There is no outcome from a real evacuation to
predict. The ladder cannot be checked against a real evacuation of this home,
and this is declared.

\paragraph{Hypothesis validity.} The ladder should behave as theory predicts on
cases with a known answer. The scorer's fixtures pin this. For example, a plan
that sends residents to a closed shelter must score below one that re-routes
them.

\paragraph{Consequential validity.} How could a system game the measure? Three
rules close the obvious routes. The ceiling is computed from what was declared,
not from the run. The clearance quantiles are computed over the ceiling, not
over survivors. And an overload invalidates only the excess residents, so no
system gains by delivering fewer.

% =====================================================================
\section{Evaluation Harness}
\label{sec:eval:harness}

One harness runs every cell of the evaluation. For each run it records a
manifest row inside the append-only log: the system, the model and its digest,
the scenario and its hash, the time scale, the replicate, the policy seed, the
git commit and whether the working tree was clean, the threshold hash, and the
platform. A run's provenance can therefore be recovered from the log itself,
not only from this document. The reproduction steps are in
Appendix~\ref{app:repro}.

\paragraph{The headless world client.} On the agentic side, a bus leg closes
only when the world acknowledges that the task is complete. At first, the only
source of that acknowledgement was the browser's vehicle animation. Cells run
without a browser therefore ran in a world where no bus ever arrived. The first
pilot measured this missing acknowledgement, not the models: all its agentic
cells scored level 0. It was declared invalid and re-run on the same plan. Every
agentic cell now runs a headless client that plays the vehicle animations from
the backend's own events. A cell refuses to start while a browser tab is
connected, because tab and client would acknowledge every leg twice.

\paragraph{A serial sweep.} The sweep cannot run in parallel on one machine.
Agent state is keyed by a fixed agent name, the agents' working directories are
fixed paths, and the route cache is a single file. It is therefore serial and
resumable. An agentic cell takes about 20 wall-minutes. The anticipated arm's
60 agentic cells took about 21 hours. The fixed order has one declared cost:
any drift in the hosted models or the hardware over time is confounded with a
run's position in the sweep.

\paragraph{Replicate blocks.} The replicate index shifts the policy seed, so
blocks are kept apart by index and by output file. Replicates 100--105 are the
pilot, 1000--1009 the confirmatory runs, and 2000 onwards the model
screenings. The scoring and summary commands read only the confirmatory file
by default. Pilot and screening runs are therefore kept out of the confirmatory
analysis mechanically, not by convention.

\paragraph{Two-pass scoring.} A cell's clean-run reference $K^\ast$ must exist
before any run in it may rise to level 4. Scoring therefore runs in two passes:
first the references, then the levels.

% =====================================================================
\section{Declared Scope}
\label{sec:eval:scope}

The protocol states in advance what it does not measure, so that an examiner
reads it here rather than finding it. Table~\ref{tab:eval:scope} condenses the
list. The limitations that arose from executing the design are discussed in
Section~\ref{sec:crit:limitations}.

\begin{table}[htbp]
\centering
\caption{What the protocol does not measure, declared before the evaluation.}
\label{tab:eval:scope}
\small
\begin{tabularx}{\textwidth}{l>{\raggedright\arraybackslash}X}
\toprule
Not measured & Why \\
\midrule
Continuous hazard exposure & No water model. The flood is one arrival time per zone, and the deadline is its discrete substitute. \\
Equity across subgroups & Registered as undefined (Section~\ref{sec:eval:metrics}). \\
Congestion and queues & No road-capacity model. \\
Per-resident trajectories & Residents are counts. The organisations are the agents, the residents are the payload. \\
Sensor noise, message delay and loss & No sensor or communication model. Adding one would test the database architecture, not the coordination policy. \\
Further ablations & Only the no-replan control. Other components are not separable, or need the declined subgroup structure. \\
Message content & The decision rubric reads structure, never prose. \\
Oracle-relative regret & No optimum is computable over an open space of novel disruptions. \\
Route-choice baselines & Route choice lives in the driver agents, which neither coordinator controls. \\
Randomised run order & The sweep is serial by construction. \\
Bit-exact reproducibility & Seeded inference is near-deterministic only. \\
The H3 capability threshold & The final set holds the strongest models, and the weaker one did not run on the novel arm. \\
The effect four novel scenarios test & It is judged on the secondary rubric, while the level is likely to tie. \\
\bottomrule
\end{tabularx}
\end{table}

The last row deserves emphasis, because it was known before the novel arm ran.
The protocol recorded that on four of the five novel scenarios, H2 would be
read on a level that was likely to tie, with the rubric's counts beside it.
How far this reached, and what it means for the answer to RQ1, is the subject
of Chapter~\ref{ch:critique}.
